\documentclass[11pt]{article}
\input{packages.tex}
\DeclareMathOperator{\sign}{sign}
\DeclareMathOperator{\Imm}{Im}

% Настройка колонтитулов
\pagestyle{fancy}
\fancyhf{} % Очистить все поля
\fancyhead[L]{Эх, Киса! Мы чужие на этом празднике жизни...}
\fancyhead[R]{\href{https://t.me/ceagest}{Автор}}
\fancyfoot[C]{\thepage}               % Справа внизу

\setlength{\headheight}{13.6pt}

\hypersetup{
    colorlinks=true,
    linkcolor=blue,
    citecolor=blue,
    filecolor=blue,
    urlcolor=blue,
}

\begin{document}

\begin{titlepage}
    \centering
    \vspace*{\fill}  % заполняет пространство сверху

    {\Huge Лекция 7. Лемма о накачке (разрастании). Теорема Майхилла-Нерода. Минимальный ПДКА.
    \par}

    \vspace*{\fill}  % заполняет пространство снизу
\end{titlepage}

\setcounter{page}{1}

\section{Лемма о накачке (разрастании)}

\begin{lemma}
    Пусть язык $\mathcal{L} \subset \Sigma^{\ast}$ регулярен. Тогда существует такое $k \in \mathbb{N}$, что для любого слова $w \in \mathcal{L}$ длины больше $k$ существуют слова $x, y, z$ такие, что $w = xyz$ и
    \begin{itemize}
        \item $\left\lvert y \right\rvert > 0$
        \item $\left\lvert xy \right\rvert \leq k$
        \item $xy^nz \in \mathcal{L}$ для любого $n \in \mathbb{N} \cup \{0\}$
    \end{itemize}
\end{lemma}

\begin{proof}
    Так как язык $\mathcal{L}$ регулярный, то он автоматный. Тогда существует ДКА $(\Sigma, Q, Q_f, \delta, q_0)$, его принимающий. По определению в этом автомате $\left\lvert Q \right\rvert$ состояний. Рассмотрим некоторое слово $w \in \mathcal{L}$ длины больше, чем $\left\lvert Q \right\rvert + 1$. Тогда, считывая автоматом такое слово, по принципу Дирихле мы побываем в какой то вершине диаграммы Мура для этого ДКА дважды, причём повтор случится не позднее, чем на считывании $\left\lvert Q \right\rvert + 1$ буквы этого слова. Получили, что в этом префиксе из $\left\lvert Q \right\rvert + 1$ букв точно есть подслово, "зацикливающее" \; некоторое состояние. При этом очевидно, что это подслово непусто. Тогда, положив $k := \left\lvert Q \right\rvert + 1$ получим то, что и требовалось.
\end{proof}

\begin{remark}
    Это лемма \textbf{не является критерием}! В самом деле, приведём пример нерегулярного языка, который удовлетворяет стандартной версии леммы:
    \[\mathcal{L} = \{a^ib^jc^k \; | \; i \neq 1, j \geq 0, k \geq 0\} \cup \{ab^ic^i \; | \; i \geq 1\}\]
\end{remark}

\begin{example}
    Рассмотрим язык
    \[\mathcal{L} = \{0^n1^n \; | \; n \in \mathbb{N}\} \subseteq \{0, 1\}^{\ast}\]
    Покажем, пользуясь леммой о накачке, что этот язык не является регулярным. \\
    От противного предположим, что $\mathcal{L}$ является регулярным. Покажем, что в таком случае выполнено отрицание леммы о накачке, доказав таким образом требуемое.
    Зафиксируем некоторое произвольное $k \in \mathbb{N}$ и рассмотрим слово 
    \[w = 0^k1^k \in \mathcal{L}\]
    Рассмотрим также некоторое произвольное разбиение $xyz$ слова $w$ такое, что
    \[\left\lvert y \right\rvert > 0, \quad \left\lvert xy \right\rvert \leq k, \quad w = 0^k1^k = xyz\]
    Так как $\left\lvert xy \right\rvert \leq k$, то подслово $xy$ рассматриваемого разбиения полностью состоит из нулей. Из этого сразу следует, что 
    \[y = 0^t, \quad t \neq 0\]
    Но тогда при $n = 2$ имеем:
    \[xy^nz = xy^2z = 0^{k + t}1^k \not \in \mathcal{L}\]
    Легко видеть, что получили то самое отрицание леммы о накачке, а значит требуемое доказано.
\end{example}

\section{Теорема Майхилла-Нерода}

\begin{definition}[Отношение Майхилла-Нерода]
    Пусть дан язык $\mathcal{L} \subset \Sigma^{\ast}$. Определим отношение Майхилла-Нерода $\equiv_{\mathcal{L}}$ на $\Sigma^{\ast}$ следующим образом:
    \[x \equiv_{\mathcal{L}} y \iff (\forall z \in \Sigma^{\ast} \hookrightarrow xz \in \mathcal{L} \iff yz \in \mathcal{L})\]
\end{definition}

\begin{theorem}[Майхилла-Нерода]
    Язык $\mathcal{L}$ регулярен тогда и только тогда, когда число классов эквивалентности по отношению Майхилла-Нерода конечно.
\end{theorem}

\begin{proof} $(\Longrightarrow): $ Пусть язык $\mathcal{L} \subset \Sigma^{\ast}$ регулярен. Тогда для этого языка существует ПДКА $(\Sigma, Q, Q_f, \delta, q_0)$, его принимающий. Пусть $q \in Q$. Заметим, что если из $q_0$ в $q$ можно попасть и по слову $x$, и по слову $y$, то $x \equiv_{\mathcal{L}} y$. В самом деле, как только мы оказались в состоянии $q$, далее нам важен лишь суффикс для двух этих слов --- либо вследствие его добавления придём в принимающее состояние, либо нет.
    Теперь вспомним, что мы взяли полный автомат, а значит считываться до конца будет любое слово $\Sigma^{\ast}$. При этом из доказанного выше следует, что слова, ведущие в одно и то же состояние, будут эквивалентны по отношению Майхилла-Нерода для $\mathcal{L}$. Тогда всего классов эквивалентности по этому отношению может быть не более $\left\lvert Q \right\rvert$. Но множество $Q$ --- конечно, а значит требуемое доказано. \\
    $(\Longleftarrow): $ Пусть язык $\mathcal{L} \subset \Sigma^{\ast}$ таков, что количество классов по отношению Майхилла-Нерода для него конечно. Так как каждый класс эквивалентности можно задать одним словом из алфавита $\Sigma$, то имеем возможность выписать это конечное количество классов в следующем виде:
    \[[\varepsilon], [w_1], \ldots, [w_k]\]
    При этом нетрудно видеть, что
    \[\forall a \in \Sigma \hookrightarrow x, y \in [w] \implies xa, ya \in [wa]\]
    Тогда, пользуясь этим чудесным фактом, мы можем построить ДКА, принимающий $\mathcal{L}$. Для этого достаточно лишь взять вершинами диаграммы Мура этого ДКА классы эквивалентности и проводить ребро с буквой $a \in \Sigma$ из $[w_i]$ в $[w_j]$, если $w_j \equiv_{\mathcal{L}} w_ia$ (состояние будет принимающим, если слово, задающее соответствующий класс эквивалентности, лежит в языке). Получили, что язык $\mathcal{L}$ автоматный, а значит регулярный, что и требовалось.
\end{proof}

\begin{remark} \label{rem:2.1}
    Детерминированный конечный автомат, построенный на классах эквивалентности в процессе доказательства теоремы выше, обладает двумя классными свойствами: он является полным и минимальным. Минимальность понимается в том смысле, что не существует ДКА с меньшим количеством состояний, который принимал бы рассматриваемый язык. \\
    Действительно, полнота очевидна, так как из любого состояния по букве определён переход в некоторый другой класс. Минимальность же следует из того, что слова, ведущие в одну и ту же вершину, эквивалентны, а значит вершин уж точно не меньше, чем классов эквивалентности. В нашем автомате их ровно столько, что и требовалось.
\end{remark}

\section{Минимальный ПДКА}

\begin{definition}
    ДКА называется \textbf{всюду определённым (полным, ПДКА)}, если:
    \[\forall a \in \Sigma, \; \forall q \in Q \hookrightarrow \delta(q, a) \neq \varnothing\]
\end{definition}

\begin{lemma}
    Для любого ДКА существует эквивалентный ПДКА.
\end{lemma}

\begin{proof}
    Доказательство состоит лишь в добавлении в рассматриваемый ДКА "мусорного" \; состояния $T$, в которое будет определён переход в новом ПДКА, если в старом ДКА он ссылался на пустое множество.
\end{proof}

\begin{lemma}
    В ПДКА с минимальным числом состояний каждое состояние соответствует одному из классов эквивалентности по отношению Майхилла-Нерода.
\end{lemma}

\begin{proof}
    Требуемое легко следует рассуждений, аналогичных рассуждениям в замечании \ref{rem:2.1}, а также из двух следующих наблюдений:
    \begin{itemize}
        \item В любом ДКА все префиксы, приводящие в одно состояние, относятся к одному классу эквивалентности.
        \item Зная классы эквивалентности, легко построить соответствующий ДКА.
    \end{itemize}
\end{proof}

\subsection{Алгоритм сведения ПДКА к минимальному ПДКА.} 
\textbf{Вход:} ПДКА $(\Sigma, Q, Q_f, \delta, q_0)$ \\
\textbf{Выход:} ПДКА $(\Sigma, Q', Q_f', \delta', q'_0)$ с минимальным числом состояний, задающий тот же язык.
\begin{enumerate}
    \item Построить начальное разбиение $\Pi = \{Q \setminus Q_f, Q_f\}$.
    \item Для каждой группы $G \in \Pi$:
    \begin{itemize}
        \item Разбить $G$ на подмножества так, чтобы состояния $u$ и $v$ были в одном подмножестве тогда и только тогда, когда для каждого $a \in \Sigma$ состояния $u$ и $v$ имеют переходы по $a$ в состояния одной и той же группы $\Pi$.
    \end{itemize}
    \item Если $\Pi_{\text{new}} = \Pi$, то $\Pi_{\text{res}} = \Pi$, иначе повторить шаг 2.
    \item $Q' = \Pi$; $q'_0 = G$, где $q_0 \in G$; $Q'_f = \{G \; | \; G \in Q', G \cap Q_f \neq \varnothing\}$; $\delta'(p', a) = q'$, если $\delta(p, a) = q$, где $p \in p'$, $q \in q'$.
    \item Удалить состояния, недостижимые из начального. Они могут появиться, если такие были в исходном.
\end{enumerate}

\begin{example}
    Рассмотрим следующий ПДКА:
    \begin{figure}[h]
    \centering
    \includegraphics[width=0.65\linewidth]{pictures/Сведение к минимальному ПДКА.jpg}
    \end{figure}
    \begin{enumerate}
        \item $\Pi = \{\{1, 2, 3, T\}, f\}$
        \begin{itemize}
            \item $\displaystyle 1 \mathop{\longrightarrow}^c \; \{1, 2, 3, T\}$
            \item $\displaystyle 2 \mathop{\longrightarrow}^c \; f$
            \item $\displaystyle 3 \mathop{\longrightarrow}^c \; f$
            \item $\displaystyle T \mathop{\longrightarrow}^c \; \{1, 2, 3, T\}$
        \end{itemize}
        \item $\Pi = \{\{1, T\}, \{2, 3\}, f\}$
        \begin{itemize}
            \item $\displaystyle 1 \mathop{\longrightarrow}^a \; \{2, 3\}$
            \item $\displaystyle T \mathop{\longrightarrow}^a \; \{1, T\}$
        \end{itemize}
        \item $\Pi = \{\{1\}, \{T\}, \{2, 3\}, f\}$
    \end{enumerate}
    Тогда полученный новый минимальный ПДКА имеет вид:
    \begin{figure}[h]
    \centering
    \includegraphics[width=0.65\linewidth]{pictures/Результат алгоритма сведения.jpg}
    \end{figure}
\end{example}

\section{Замкнутость регулярных языков}

Пусть $\mathcal{A}$ и $\mathcal{B}$ --- регулярные языки над алфавитом $\Sigma$. Тогда регулярными также являются:

\begin{itemize}
    \item $\mathcal{A} \cup \mathcal{B}, \; \mathcal{A} \cdot \mathcal{B}, \; \mathcal{A}^{\ast}$ (по определению)
    \item $\overline{\mathcal{A}} = \Sigma^{\ast} \setminus \mathcal{A}$
    \item $\mathcal{A} \cap \mathcal{B}$
    \item $\mathcal{A} \setminus \mathcal{B}$
    \item гомоморфизм, реверс и др.
\end{itemize}

\section{Декартово произведение ДКА}

\begin{definition}
    \textbf{Прямым произведением} двух ДКА 
    \[A_1 = (\Sigma_1, Q_1, Q_{f1}, \delta_1, q_{01}) \quad \text{ и } \quad A_2 = (\Sigma_2, Q_2, Q_{f2}, \delta_2, q_{02})\]
    называется ДКА $(\Sigma, Q, Q_f, \delta, q_0)$, где:
    \begin{itemize}
        \item $\Sigma = \Sigma_1 \times \Sigma_2$
        \item $Q = Q_1 \times Q_2$
        \item $Q_f = Q_{f1} \times Q_{f2}$
        \item $\delta(\left\langle q_1, q_2 \right\rangle, a) = \left\langle \delta_1 (q_1, a), \delta_2 (q_2, a)\right\rangle$
        \item $q_0 = \left\langle q_{01}, q_{02} \right\rangle$
    \end{itemize}
\end{definition}

\begin{lemma}
    Прямое произведение ДКА задаёт пересечение языков, задаваемых исходными.
\end{lemma}

\end{document}